% Pursuing Stacks (cotes 134-2 to 134-8): archive cote 134-7, site page 32 (PDF page 33), typescript page 474.
% typescript. Draft compared: ps6.tex:1488-1560.
% AI-transcribed by Claude Opus 5.5 (claude-opus-5-5, adaptive thinking, effort high), Message Batches API, batch ps3 (msgbatch_01XdFmuVVBAcgc7EFxdtCqsj), 2026-11-12; reviewed against both supplied scan halves; see CHECK_DECISIONS.tsv.
\scanpage{134-7/32}
\mspage{474}

The description I just recalled of U(k) is the one most intuitive to
my mind,  other people may prefer the more computational one given on
page 451 in terms of $\hand{\Gamma}_k^{\hand{\wedge}}\,(M)$  (endowed with it's augmentation to k and it's
diagonal map),  which is of importance in it's own right.  It shows
the existence of a canonical k-\uline{linearization} \uline{functor}
\[ \text{(J)} \qquad L : U(k) \hand{\longrightarrow} Ab_k \quad , \]
giving rise to the commutative diagram (up to can. isom.)
\deleted{\xout{$Ab_k$x\allowbreak{}x\allowbreak{}x\allowbreak{}x\allowbreak{}x\allowbreak{}x\allowbreak{}$W$x\allowbreak{}x\allowbreak{}x\allowbreak{}x\allowbreak{}x\allowbreak{}x\allowbreak{}x\allowbreak{}$Ab_k$x\allowbreak{}x\allowbreak{}W(k)}}
% Hand-drawn arrows are marked blue; typed objects remain plain.
\diagram{
\text{(J')}\qquad
\begin{tikzcd}
Ab_k \arrow[color=blue,rr, "W"] \arrow[color=blue,dr, "\hand{\Gamma}_k^{\hand{\wedge}}"'] \& \& U(k) \arrow[color=blue,dl, "L"] \\
\& Ab_k \&
\end{tikzcd}
\quad \xout{,}
}
where
\[ \hand{\Gamma}_k^{\hand{\wedge}}\,(M) = \hand{\prod}_{i\hand{\geq}0} \hand{\Gamma}_k^{\,i}(M) \qquad . \]
This linearization is not quite compatible with ring extension,  it beco-%
mes so only when we view it as a functor with values,  not just in the
category $Ab_k$ of k-modules,  but of separated and complete \xout{k-modules for}
linearly topologized k-modules,  the ring extension functor for these
being \uline{completed} tensor product.  This is a little (or big ?) technical
\margin{x} drawback for this notion of linearization.  We have a canonical embedding
\[ \text{(K)} \qquad \hand{x \mapsto \exp(x) :} \qquad X \hand{\longrightarrow} \xout{\unsure{W(L(X))}xxx}\ W^{\hand{\wedge}}\,(L(X)) \qquad \hand{\text{(cf. p.450, (n))}} \]
functorial in X,  \xout{which satisfies a} where for a topological k-module M
as above,  described as a filtering inverse limit of discrete ones $M_i$ ,
we define
\[ \text{(K')} \qquad \xout{W}\ W^{\hand{\wedge}}\,(M) = (k' \hand{\longmapsto} \underset{\hand{\leftarrow\, i}}{\underline{\lim}}\ M_{\unsure{i}} \otimes_k k' \ ) \quad . \]
The map (K) has a universal property with respect to all possible maps
$X \hand{\to} W^{\hand{\wedge}}\,(M)$  with M a linearly topologized separated and complete k-module,
which acounts for it's role as ``linearization''.  It should be noted here
that the map (C) on page (467) does'nt quite exist,  we have corrected
this point here -  definitely we cannot in (K) replace $W^{\hand{\wedge}}$ by W .  Of cour-%
se linearization L (or it's variant $L_{pt}$)  does'nt commute in any sense
whatever to restriction of ground ring.

The image of  X  in $W^{\hand{\wedge}}\,(L(X))$  is caracterized by the simple
formulae (q) p. 451 .  Maps from X to Y  may be described as just continu-%
ous \inserted[typed]{k-linear} maps from  L(X) to L(Y),  compatible with augmentations and diag\unsure{o}nal
maps.

\margin{32}

% Visual collation: both supplied image halves inspected; see CHECK_DECISIONS_134_7_8.tsv.
